\subsection{Prerequisite lemmas}


We require some lemmas before proving Theorem~\ref{intro_version_p}, and we introduce more notation to do so.
For a complex number \(\alpha\), let \(\bar{\alpha}\) denote the complex conjugate of \(\alpha\), and let \(\| \alpha \| = \sqrt{\alpha \bar{\alpha}} \in [0,\infty)\) denote the usual norm of \(\alpha\).
Define the function \(D : \N \to \N \cup \{0\}\) by
\begin{equation}
D(n) =
\begin{cases}
0 & \text{ if } n = 1, \\
\max \{ s \in \N : s \mid n \text{ and } s < n\} & \text{ if } n > 1.
\end{cases}
\end{equation}
In other words, \(D(n)\) is the largest proper divisor of \(n\), unless \(n = 1\), and \(D(1) = 0\).
Note that \(D(n) \leq n/2\) for all \(n \in \N\).
In this section, we also make use of \myemph{big \(O\) notation}.
Recall that if \(S\) is an infinite subset of \(\N\) and \(f,g : S \to [0,\infty)\), then we write \(f = O(g)\) to denote that there exist \(m \in [0,\infty)\) and \(s_0 \in S\) such that \(f(s) \leq m \cdot g(s)\) for all \(s > s_0\).
In other words, \(f\) is big \(O\) of \(g\) if some constant multiple of \(g(s)\) is an upper bound on \(f(s)\) for all sufficiently large \(s \in S\).

\begin{remark}
	Our aim is to use big \(O\) statements to evaluate limits as \(q \to \infty\).
	Therefore, in all expressions involving big \(O\) notation in this section, we will take \(S\) to be the set of prime powers and use \(q\) to denote the argument of each relevant function.
\end{remark}


\begin{lemma}
\label{degree_limit_lemma}
For all \(n \in \N, \, d \mid n\), and \(r \in \{0,\ldots,\tfrac{n}{d}-1\}\), we have
\begin{equation}
\frac{1}{\deg_{n,d,r} (q)}
=
O \left (
q^{-e(n,d,r)}
\right ),
\end{equation}
where we define
\begin{equation}
\label{definition_of_e}
e(n,d,r)
=
d\binom{r+1}{2} + \binom{n+1}{2} - d\binom{\tfrac{n}{d} + 1}{2} + dr\left ( \tfrac{n}{d}-1-r \right ).
\end{equation}
Moreover, \(e(n,d,r)\) is positive unless \(d = 1\) and \(r = 0\), and \(e(n,1,0) = 0\).
\end{lemma}

\begin{proof}
The first claim follows from the fact that \(\deg_{n,d,r}\) is a rational function of~\(q\) and then checking the degrees in \(q\) of the various polynomials which comprise \(\deg_{n,d,r} (q)\).
To prove the second claim, first observe that \(e(n,d,r)\) is a quadratic polynomial in \(r\) with critical point \(r = \tfrac{n}{d} - \tfrac{1}{2}\) and leading coefficient \(-d/2\).
This implies \(e(n,d,r)\) is increasing for \(r \in \{0,\ldots,\tfrac{n}{d}-1\}\).
The minimum value of \(e(n,d,r)\) on \(\{0,\ldots,\tfrac{n}{d}-1\}\) is therefore achieved at \(r = 0\).
Observe that \(e(n,d,0) = \tfrac{n}{2}\left ( n - \tfrac{n}{d}\right )\), which is positive unless \(d = 1\).
Therefore, \(e(n,d,r) > 0\) if \(d > 1\).
In the case \(d = 1\), we have \(e(n,1,r) = -\tfrac{1}{2}r^2 + (n-\tfrac{1}{2})r\), which is positive unless \(r = 0\).
The result follows.
\end{proof}







\begin{lemma}
\label{helvetica_scenario}
For all \(n \in \N, d \mid n, r \in \{0,\ldots,n/d-1\}\), we have
\begin{equation}
\max_{\substack{f \in \CF_d (q) \\ f(z) \neq z-1}}
\frac{\left \| \sum_{g \in \CT_{(n)}(q)} \chi^{f,r}(g) \right \|}{\# \CT_{(n)}(q)}
=
O \left ( q^{D(n)-n} \right ).
\end{equation}
\end{lemma}

\begin{proof}
Applying Theorem~\ref{main_lemma_in_intro}, Corollary \ref{MNrule_ncycle}, \eqref{size_of_Fm}, Corollary \ref{rss_cc_and_cycle_type_sizes}, and Lemma \ref{pre_pre_main}, we have
\begin{equation}
\label{final_row}
\frac{
1
}{
\# \CT_{(n)} (q)
}
\left \|
\sum_{g \in \CT_{(n)} (q)} \chi^{f,r} (g)
\right \|
=
\frac{d \cdot \left | \sum_{s \mid n} \moebius(n/s)(q^s-1) \delta_{q^s - 1 \mid \ell_{f} [n/d]_{q^d}} \right | }{\sum_{s \mid n} \moebius(n/s)(q^s-1)}
\end{equation}
for all prime powers \(q\) and \(f \in \CF_d (q)\).
The denominator on the right side of \eqref{final_row} is a degree-\(n\) polynomial in \(q\), independent of \(f\).
However, the numerator on the right side of \eqref{final_row} is not necessarily a polynomial in \(q\) at all, as it also depends on \(\ell_{f}\) which can vary with \(q\), and the sum is inside of an absolute value.
Fortunately, if \(q^n - 1\) does not divide \(\ell_{f} [n/d]_{q^d}\), then
\begin{equation}
\label{D_n_bound}
\begin{aligned}
\left | \sum_{s \mid n} \moebius(n/s) (q^s-1) \delta_{q^s-1 \mid \ell_f \cdot [n/d]_{q^d}} \right |
& \leq
\sum_{\substack{s \mid n \\ s < n}}
\left | \moebius (n/s) (q^s-1) \right | \\
& \leq
\sum_{\substack{s \mid n \\ s < n}}
q^s
<
1 + \sum_{\substack{s \mid n \\ s < n}}
q^s.
\end{aligned}
\end{equation}
Therefore, in this case
\begin{equation}
\label{final_row_2}
\frac{
1
}{
\# \CT_{(n)} (q)
}
\left \|
\sum_{g \in \CT_{(n)} (q)} \chi^{f,r} (g)
\right \|
\leq
d \cdot \frac{1 + \sum_{s \mid n, \, s< n} q^s }{\sum_{s \mid n} \moebius(n/s)(q^s-1)}.
\end{equation}
Observe that \(1 + \sum_{s \mid n, \, s < n} q^s\) is a degree-\(D(n)\) polynomial in \(q\), and the right side of \eqref{final_row_2} is independent of \(f\).
Moreover, the condition \(q^n-1 \nmid \ell_f \cdot [n/d]_{q^d}\) is equivalent to \(f(z) \neq z-1\) because
\[
q^n-1 \mid \ell_f \cdot [n/d]_{q^d}
\iff q^d-1 \mid \ell_f
\iff f(1) = 0
\iff f(z) = z-1.
\]
The result now follows from computing the maximum of \eqref{final_row_2} over \mbox{\(f \in \CF_d(q) \setminus \{z-1\}\).}
\end{proof}









\begin{lemma}
\label{char_sum_norm_lemma}
For all \(n \in \N, \, \mu \vdash n, \, d \mid n\), \(r \in \{0,\ldots,\tfrac{n}{d}-1\}\), we have
\begin{equation}
\max_{f \in \CF_d(q)}
\frac{\left \| \sum_{g \in \CT^\Box_\mu(q)} \chi^{f,r}(g) \right \|}{\gamma_n(q)}
=
O(1).
\end{equation}
\end{lemma}

\begin{proof}
Consider a fixed prime power \(q\) and polynomial \(f \in \CF_d (q)\) to begin.
By Theorem~\ref{main_lemma_in_intro}, if some part of \(\mu\) is not divisible by \(d\), then \(\sum_{g \in \CT_\mu^\Box (q)} \chi^{f,r}(g) = 0\), which satisfies the claim.
So assume there exists \(\tilde{\mu} \vdash n/d\) such that \(\mu = d \tilde{\mu}\).
Recall that, by Theorem~\ref{fulmans_characterization}, the conjugacy classes in \(\CT_\mu^\Box (q)\) are in bijection with subsets \(\{h_1,\ldots,h_{\ell(\mu)}\} \subset \CF (q)\) of distinct polynomials with \(\deg h_i = \mu_i\) for each \(i \in \{1,\ldots,\ell(\mu)\}\).
By Theorem~\ref{main_lemma_in_intro}, Corollary~\ref{rss_cc_and_cycle_type_sizes}, and the fact that characters are constant on conjugacy classes, we have that
\(
\sum_{g \in \CT_\mu^\Box (q)}
\chi^{f,r} (g)\)
equals
\begin{equation}
\label{everybodys_gotta_learn_sometime}
\frac{\gamma_n(q) (-1)^{\tfrac{n}{d}(d-1)} \chi^{(n/d-r,1^r)}_{\tilde{\mu}}}{\prod_{i=1}^{\ell(\mu)} (q^{\mu_i}-1)}
\sum_{ \substack{\{ h_1,\ldots,h_{\ell(\mu)}\} \subset \CF (q) \\ \deg h_i = \mu_i \forall i} }
\prod_{i=1}^{\ell(\mu)}
\frac{1}{\tilde{\mu}_i}
\sum_{\substack{\alpha_i \in \F_{q^{\mu_i}} \\ h_i(\alpha_i) = 0}}
\te(\alpha_i)^{\ell_f \cdot [\tilde{\mu}_i]_{q^d}}.
\end{equation}
We can now separate the outer sum in \eqref{everybodys_gotta_learn_sometime} according to the degrees of the distinct polynomials \(h_i \in \CF_{\mu_i} (q)\).
Doing so transforms \eqref{everybodys_gotta_learn_sometime} into
\begin{equation}
\frac{\gamma_n(q) (-1)^{\tfrac{n}{d}(d-1)} \chi^{(n/d-r,1^r)}_{\tilde{\mu}}}{\prod_{i=1}^{\ell(\mu)} (q^{\mu_i}-1)}
\prod_{s \geq 1}
\left (
\frac{d}{s}
\right )^{m_s(\mu)}
\hspace{-0.4cm} \sum_{ \{p_1,\ldots,p_{m_s(\mu)}\} \subset \CF_s (q) }
\prod_{i=1}^{m_s(\mu)}
\sum_{\substack{\beta_i \in \F_{q^s} \\ p_i(\beta_i)=0}}
\te(\beta_i)^{\ell_f \cdot [s/d]_{q^d}}.
\end{equation}
Computing the norm, applying the triangle inequality, recalling that \(\te\) maps into the unit circle in \(\Cbb\), and applying Corollary~\ref{rss_cc_and_cycle_type_sizes} gives
\begin{align*}
& \left \|
\sum_{g \in \CT^\Box_\mu (q)} \chi^{f,r} (g)
\right \| \\
\leq &
\frac{\gamma_n(q) \left | \chi^{(n/d-r,1^r)}_{\tilde{\mu}} \right |}{\prod_{i=1}^{\ell(\mu)} (q^{\mu_i}-1)}
\prod_{s \geq 1}
\left (
\frac{d}{s}
\right )^{m_s(\mu)}
\sum_{ \{p_1,\ldots,p_{m_s(\mu)}\} \subset \CF_s (q) }
\prod_{i=1}^{m_s(\mu)}
\sum_{\substack{\beta_i \in \F_{q^s} \\ p_i(\beta_i)=0}}
\left \|
\te(\beta_i)^{\ell_f \cdot [s/d]_{q^d}}
\right \| \\
= &
\frac{\gamma_n(q) \left | \chi^{(n/d-r,1^r)}_{\tilde{\mu}} \right |}{\prod_{i=1}^{\ell(\mu)} (q^{\mu_i}-1)}
\prod_{s \geq 1}
\left (
\frac{d}{s}
\right )^{m_s(\mu)}
\sum_{ \{p_1,\ldots,p_{m_s(\mu)}\} \subset \CF_s (q) }
s^{m_s(\mu)} \\
= &
\frac{\gamma_n(q) \left |\chi^{(n/d-r,1^r)}_{\tilde{\mu}} \right |}{\prod_{i=1}^{\ell(\mu)} (q^{\mu_i}-1)}
d^{\sum_{s \geq 1} m_s(\mu)}
\prod_{s \geq 1}
\binom{\# \CF_s (q)}{m_s(\mu)} = \# \CT_\mu^\Box (q) \cdot \left | \chi^{(n/d-r,1^r)}_{\tilde{\mu}} \right | \cdot d^{\, \ell(\mu)}.
\end{align*}
Thus,
\begin{equation}
\label{constant_bound}
\frac{1}{\gamma_n(q)}
\left \|
\sum_{g \in \CT_\mu^\Box (q)} \chi^{f,r} (g)
\right \|
\leq
\frac{\# \CT_\mu^\Box (q)}{
\gamma_n (q)}
\cdot \left | \chi^{(n/d-r,1^r)}_{\tilde{\mu}} \right | \cdot d^{\, \ell(\mu)}.
\end{equation}
The right side of \eqref{constant_bound} does not depend on \(f\), which implies
\begin{equation}
\label{constant_bound_2}
\max_{f \in \CF_d (q)}
\frac{1}{\gamma_n(q)}
\left \|
\sum_{g \in \CT_\mu^\Box (q)} \chi^{f,r} (g)
\right \|
\leq
\frac{\# \CT_\mu^\Box (q)}{
\gamma_n (q)}
\cdot \left | \chi^{(n/d-r,1^r)}_{\tilde{\mu}} \right | \cdot d^{\, \ell(\mu)}.
\end{equation}
Moreover, by Corollary~\ref{dense}, for sufficiently large \(q\), the right side of \eqref{constant_bound_2} is arbitrarily close to the constant value \(|\chi_{\tilde{\mu}}^{(n/d-r,1^r)}| \cdot d^{\ell(\mu)} / z_\mu\).
The result follows.
\end{proof}















\subsection{Proof of asymptotic result}

Recall that we have defined
\begin{equation}
p_{k,\mu} (q) = \frac{g_{k,\mu}(q)}{\# \CT_{(n)} (q)^k}
\quad \text{and} \quad
p_{k,\mu}^\Box (q) = \frac{g_{k,\mu}^\Box(q)}{\# \CT_{(n)} (q)^k}
\end{equation}
for all \(n,k \in \N\), \(k \geq 2\), and \(\mu \vdash n\).
We can now prove our asymptotic result, Theorem~\ref{intro_version_p}, which states
\begin{equation}
\lim_{q \to \infty}
p_{k,\mu} (q)
=
\lim_{q \to \infty}
p_{k,\mu}^\Box (q)
=
\frac{1}{z_\mu}.
\end{equation}


\begin{proof}[Proof~of~Theorem~\ref{intro_version_p}]
We will first prove that \(\lim_{q \to \infty} p_{k,\mu}^\Box (q) = 1 / z_\mu\).
From this, it follows that \(\lim_{q \to \infty} p_{k,\mu} (q) = 1 / z_\mu\) because, for all prime powers \(q\), we have \(p_{k,\nu} (q) \geq p_{k,\nu}^\Box (q)\) for all \(\nu \vdash n\) and \(\sum_{\nu \vdash n} p_{k,\nu} (q) = 1 = \sum_{\nu \vdash n} 1/z_\nu\).

Consider the following formulation of \(p_{k,\mu}^\Box (q)\).
By its definition \eqref{pBox_defs} and by Corollary~\ref{simpler_ff_in_situ_corollary}, we have
\begin{equation}
\label{pBox_formulation}
p_{k,\mu}^\Box (q) =
\sum_{d \mid n}
\sum_{r = 0}^{\tfrac{n}{d} - 1}
\sum_{f \in \CF_d (q)}
\left (
\frac{\sum_{g \in \CT_{(n)}(q)} \chi^{f,r}(g)}{\# \CT_{(n)}(q)}
\right )^k
\left (
\frac{\sum_{h \in \CT^\Box_\mu(q)} \chi^{f,r}(h)}{\gamma_n(q) \cdot \deg_{n,d,r}(q)^{k-1}}
\right ).
\end{equation}
We want to compute \(\lim_{q \to \infty} p_{k,\mu}^\Box (q)\), but the index set for the summation over \(f \in \CF_d (q)\) in \eqref{pBox_formulation} itself depends on \(q\).
Therefore, for each \(d \mid n\), \(r \in \{0,\ldots,n/d-1\}\), and \(f \in \CF_d (q)\) we define
\begin{equation}
\Phi_{k,\mu,d,r} (q)
=
\sum_{f \in \CF_d (q)}
\left (
\frac{\sum_{g \in \CT_{(n)}(q)} \chi^{f,r}(g)}{\# \CT_{(n)}(q)}
\right )^k
\left (
\frac{\sum_{h \in \CT^\Box_\mu(q)} \chi^{f,r}(h)}{\gamma_n(q) \cdot \deg_{n,d,r}(q)^{k-1}}
\right )
\end{equation}
so that
\begin{equation}
p^\Box_{k,\mu} (q)
=
\sum_{d \mid n}
\sum_{r = 0}^{\tfrac{n}{d}-1}
\Phi_{k,\mu,d,r} (q),
\end{equation}
where the number of terms in the summation is fixed, even as \(q\) varies.
Theorem~\ref{intro_version_p} now follows from Lemma~\ref{Phi_behavior} below, which computes the limiting behavior of \(\Phi_{k,\mu,d,r} (q)\) for each \(d \mid n\) and \(r \in \{0,\ldots,n/d-1\}\).
\end{proof}

\begin{lemma}
\label{Phi_behavior}
For all \(n,k \in \N, \mu \vdash n, d \mid n\), and \(r \in \{0,\ldots,n/d-1\}\), we have
\begin{equation}
\lim_{q \to \infty}
\Phi_{k,\mu,d,r} (q)
=
\begin{cases}
0         & \text{ if } d > 1 \text{ or }  r > 0, \\
1 / z_\mu & \text{ if } d = 1 \text{ and } r = 0.
\end{cases}
\end{equation}
\end{lemma}

\begin{proof}
Consider first the case that \(d > 1\) and \(r\) is arbitrary.
Observe that
\(
\|
\Phi_{k,\mu,d,r} (q)
\|
\)
is bounded above by
\begin{equation}
\label{naive_max_bound}
\frac{
\# \CF_d (q)
}{
\deg_{n,d,r} (q)^{k-1}
}
\cdot
\max_{f \in \CF_d (q)}
\left (
\frac{\left \| \sum_{g \in \CT_{(n)}(q)} \chi^{f,r}(g) \right \|}{\# \CT_{(n)}(q)}
\right )^k
\cdot
\max_{f \in \CF_d(q)}
\frac{\left \| \sum_{h \in \CT^\Box_\mu(q)} \chi^{f,r}(h) \right \|}{\gamma_n(q)}
\end{equation}
for all prime powers \(q\).
We proceed to investigate the asymptotic dependence on \(q\) of \eqref{naive_max_bound}.
Recall from \eqref{size_of_Fm} that \(\# \CF_d (q) = O(q^d)\).
Combining this with Lemmas \ref{degree_limit_lemma},
\ref{helvetica_scenario},
and
\ref{char_sum_norm_lemma},
we have
\begin{equation}
\| \Phi_{k,\mu,d,r} (q) \|
=
O \left ( q^{d + k (D(n) - n) - (k-1) \cdot e(n,d,r)} \right ).
\end{equation}
By hypothesis, \(k \geq 2\), implying \(d  + k \cdot (D(n) - n) \leq d - k \cdot n/2 \leq 0\).
Moreover, by Lemma~\ref{degree_limit_lemma}, \((k-1)\cdot e(n,d,r) > 0\).
It follows that \(\lim_{q \to \infty} \Phi_{k,\mu,d,r} (q) = 0\) if \(d > 1\).

Next, consider the case \(d = 1\).
Observing that \(z-1 \in \CF_1 (q)\) for all prime powers \(q\) and applying Theorem~\ref{sn_stuff}, we can rewrite \(\Phi_{k,\mu,1,r} (q)\) as
\begin{align}
\Phi_{k,\mu,1,r} (q)
& =
\frac{\# \CT_\mu^\Box}{\gamma_n (q)}
\cdot
\frac{(-1)^{rk} \chi^{(n-r,1^r)}_\mu}{\left ( q^{\binom{r+1}{2}} {\genfrac[]{0pt}{1}{n - 1}{r}}_{q} \right )^{k-1}}
\label{the_z_minus_1_term}
\\
& + \label{not_z_minus_1_term}
\sum_{\substack{f \in \CF_1 (q) \\ f(z) \neq z-1}}
\left (
\frac{\sum_{g \in \CT_{(n)}(q)} \chi^{f,r}(g)}{\# \CT_{(n)}(q)}
\right )^k
\left (
\frac{\sum_{h \in \CT^\Box_\mu(q)} \chi^{f,r}(h)}{\gamma_n(q) \cdot \deg_{n,1,r}(q)^{k-1}}
\right ).
\end{align}
We repeat the same analysis as before, but apply it only to \eqref{not_z_minus_1_term}.
Observe that \eqref{not_z_minus_1_term} is bounded above by
\begin{equation}
\label{not_z_minus_1_term_bound}
\frac{
(\# \CF_1(q)) - 1
}{
\deg_{n,1,r} (q)^{k-1}
}
\cdot
\max_{\substack{f \in \CF_1 (q) \\ f(z) \neq z-1}}
\left (
\frac{\left \| \sum_{g \in \CT_{(n)}(q)} \chi^{f,r}(g) \right \|}{\# \CT_{(n)}(q)}
\right )^k
\cdot
\max_{\substack{f \in \CF_1 (q) \\ f(z) \neq z-1}}
\frac{\left \| \sum_{h \in \CT^\Box_\mu(q)} \chi^{f,r}(h) \right \|}{\gamma_n(q)}
\end{equation}
Applying Lemmas \ref{degree_limit_lemma}, \ref{helvetica_scenario}, and \ref{char_sum_norm_lemma} again, we see that \eqref{not_z_minus_1_term_bound} is
\begin{equation}
\label{d_1_non_z_minus_1}
O \left (
q^{1 + k (D(n) - n) - (k-1) \cdot e(n,1,r)}
\right ).
\end{equation}
Observe that \(1 + k(D(n) - n) < 0\) even if \(n = 1\) due to the fact that \(k \geq 2\) and \(D(1)=0\), and \((k-1) \cdot e(n,1,r) \geq 0\) by Lemma~\ref{degree_limit_lemma}.
Therefore, \(1 + k \cdot (D(n) - n) - (k-1) \cdot e(n,1,r)\) is negative.
It follows that the limit as \(q \to \infty\) of \eqref{not_z_minus_1_term} is zero.
By Corollary~\ref{dense}, the limit as \(q \to \infty\) of \eqref{the_z_minus_1_term} equals
\begin{equation}
\begin{cases}
0 & \text{ if } r > 0, \\
1/z_\mu & \text{ if } r = 0.
\end{cases}
\end{equation}
The result now follows, as \eqref{the_z_minus_1_term} makes the only potentially nontrivial contribution in the limit \(q \to \infty\).
\end{proof}










